\section{Conclusion}
We presented a compact causal temporal convolutional network that reads only proprioception from a Unitree Go2 and answers three questions at once: whether a leg is entangled, which of the four legs (FR/FL/RR/RL) is affected, and how severe the entanglement is on a $[0,1]$ scale. The model is small (roughly $136$k parameters, about $0.5$~MB) and its strictly left-padded convolutions use only past and present samples, so a single forward pass runs in about $1$~ms per window on a single CPU thread, under the $2$~ms budget at $500$~Hz. Evaluation was leakage-safe by grouping whole recordings across the train/validation/test split, and leave-one-recording-out cross-validation over $15$ folds gave $\mathrm{F1}=0.807\pm0.142$ (Figure~\ref{fig:loro}); on the held-out test split the detector reached $\mathrm{F1}=0.809$ and $\mathrm{ROC\text{-}AUC}=0.956$, far ahead of the thigh-torque rule-based baseline at $\mathrm{F1}=0.222$ under an identical protocol (Table~\ref{tab:detection}). A targeted, data-centric augmentation step suppressed the field false alarms we observed on stationary and backward-walking segments while recovering the under-represented front-right event (Figure~\ref{fig:field}). Two limitations bound these claims: every detection metric is offline, computed on recorded-CSV replay rather than live on-robot operation, and the reported severity is a physics-grounded, calibrated index rather than a measured contact force, since no ground-truth severity labels exist.